\documentclass[11pt]{article}
\usepackage{amsmath,amssymb,amsthm}
\usepackage[margin=1in]{geometry}
\usepackage{hyperref}

\newtheorem{theorem}{Theorem}
\newtheorem{lemma}{Lemma}
\newtheorem{corollary}{Corollary}
\newtheorem{remark}{Remark}

\newcommand{\Term}{\mathrm{Term}}
\newcommand{\Deriv}{\mathsf{Deriv}}
\newcommand{\thmT}{\mathsf{thmT}}
\newcommand{\evalU}{\mathsf{evalU}}
\newcommand{\parse}{\mathsf{parse}}
\newcommand{\Kcode}{\mathsf{Kcode}}
\newcommand{\outK}{\mathsf{outKdef}}
\newcommand{\closeW}{\mathsf{closeW}}
\newcommand{\codeFalse}{\mathsf{codeFalse}}
\newcommand{\codeF}{\mathsf{code}}
\newcommand{\num}{\mathsf{num}}
\newcommand{\Fun}[1]{\mathsf{Fun}_{#1}}
\newcommand{\cgFun}{\mathsf{cgFun}}
\newcommand{\cgFalse}{\mathsf{cgFalse}}
\newcommand{\CGIcore}{\mathsf{CGI\_core\_num\_raw}}
\newcommand{\CGIcoreself}{\mathsf{CGI\_core\_num\_raw\_self}}
\newcommand{\Construct}{\mathsf{Construct}}
\newcommand{\Discharge}{\mathsf{DischargeKdef}}
\newcommand{\Chain}{\mathsf{ChainKdef}}
\newcommand{\Lstar}{L^{\!\star}}
\newcommand{\kmax}{k_{\max}}
\newcommand{\Sig}{\Sigma}

\title{Chaitin--G\"odel I in BRA4:\\
  a uniform-in-$w$ contradiction map}
\author{}
\date{}

\begin{document}
\maketitle

\section{Chaitin's argument, recalled}

Chaitin's proof of G\"odel's first incompleteness theorem~\cite{Chaitin1971,Chaitin1974}
runs as follows. Fix a sound theory $T$ that interprets enough arithmetic, and fix a
universal partial computable interpreter $U$ with a Kolmogorov-complexity bound
$K_U(x) := \min\{|p| : U(p)\downarrow = x\}$. There is a constant $c$
(depending only on $T$ and $U$) such that, for every $\Lstar$ and every
integer $x'$,
\[
  T \vdash K_U(\underline{x'}) > \Lstar
  \quad\Longrightarrow\quad
  \text{$T$ is inconsistent, provided } \Lstar > c + \log \Lstar.
\]
The contradiction is exhibited by a single short program $g_{\Lstar}$:
enumerate $\thmT(0), \thmT(1), \dots$ until a proof of the form
$\codeF{K_U(\underline{?}) > \Lstar}$ is found, then output the
witnessed integer~$x'$. The size of $g_{\Lstar}$ is $c + \log \Lstar$,
strictly below $\Lstar$, so the very witness $x'$ that $T$ proves to be
$K_U$-incompressible is itself $g_{\Lstar}$-output --- a contradiction
to the size lemma internal to $T$.

\section{The BRA4 formalisation: three layers}

\subsection{The core: \texorpdfstring{$\CGIcore$}{CGI core num-raw}}

The shipped core theorem (commit \texttt{2fce05d},
\texttt{BRA4/ChaitinG1CoreNumRaw.agda}) is:
\begin{align*}
  \CGIcore :{}
    & (w\ x : \Term) \to {} \\
    & \Deriv\,(\mathsf{eqF}\,(\thmT\,w)\,(\Kcode_{\Lstar}\,x)) \to {} \\
    & \Sig\,\Term\,(\lambda z.\, \Deriv\,(\mathsf{eqF}\,(\thmT\,z)\,\codeFalse)).
\end{align*}
No closure hypothesis on $w$, no \textsf{isNat}, no $\mathrm{Con}(T)$. The
``small program'' $g_{\Lstar}$ is realised by the closed BRA function
\texttt{gLcodeDef}~$\Lstar$; the search by \texttt{FirstHit.Search.gRec};
the universal interpreter by \texttt{evalU}; the chain of operational
steps closing $\evalU$ on \texttt{gLcodeDef}'s mu-loop by
\texttt{ChainKdef.dEval\_witness}; the size atom by
\texttt{dLenStarDef}. The output $z(w)$ is
\[
  \mathsf{cmp}\,
    (\mathsf{cmp}\,(\mathsf{exfProof}\,(\mathsf{D}\,g_L\,n_0\,x')\,\codeFalse)\,
                    \mathsf{cPos})\,
    (\mathsf{cmp}\,\mathsf{outerWrap}\,\mathsf{cSizeProofDef}),
\]
i.e.\ a literal encoded modus-ponens / ex-falso assembly.

\subsection{The $x$-free form: \texorpdfstring{$\CGIcoreself$}{CGI core num-raw self}}

Whether $\thmT\,w$ has the form $\Kcode_{\Lstar}\,(\_)$ is read off by
the BRA function $\outK_{\Lstar} = \mathsf{decode}\,\circ\,
\mathsf{projKdef}\,\circ\,\thmT$. Hence the parameter $x$ in the
hypothesis is redundant: it can only be $\outK_{\Lstar}\,w$. The
self-referential form, shipped in
\texttt{BRA4/ChaitinG1CoreNumRawSelf.agda}, is
\begin{align*}
  \CGIcoreself :{}
    & (w : \Term) \to {} \\
    & \Deriv\,(\mathsf{eqF}\,(\thmT\,w)\,(\Kcode_{\Lstar}\,(\outK_{\Lstar}\,w))) \to {} \\
    & \Sig\,\Term\,(\lambda z.\, \Deriv\,(\mathsf{eqF}\,(\thmT\,z)\,\codeFalse)),
\end{align*}
and the proof is one line:
$\CGIcoreself\,w\,h := \CGIcore\,w\,(\outK_{\Lstar}\,w)\,h$.

\subsection{The decomposition: $\cgFun$ and $\cgFalse$}

The $\Sigma$-payload above naturally factors as
\[
  \cgFun_{*} : (w : \Term) \to \Deriv(\dots) \to \Term,
  \qquad
  \cgFalse_{*} : (w : \Term)\,(d : \dots) \to \Deriv(\mathsf{eqF}\,(\thmT\,(\cgFun_{*}\,w\,d))\,\codeFalse),
\]
the two field projections of $\Sigma$. The first observation is that
the produced \emph{Term} is in fact $d$-independent: every Term passing
through \texttt{DischargeKdef} / \texttt{ChainKdef} /
\texttt{cgiClash} is built from $\closeW\,w$ and a finite set of
constants (the BRA combinators \texttt{gRec}, \texttt{predFlipDef}~$\Lstar$,
\texttt{outKdef}~$\Lstar$, \texttt{runProg}, \texttt{fuelMu\_fun},
\texttt{fuelG}, \texttt{thm12\_Fun2}~\texttt{runProg}, and the
tag-/pi-/sb-encoding constants). This is because
$\mathsf{leastNumber}\,B\,\textit{witness}
   = \mathsf{mkLeast}\,(\mathsf{g}\,(s\,B))\,\dots$
discards the witness for its first component (the Term $\kmax$),
and the witness Derivs are only used to construct the proof half
\textit{isPf}, not the Term \textit{pf}.

The shipped file \texttt{BRA4/CgFun.agda} therefore exposes a uniform-in-$w$
meta-Agda function
\[
  \cgFun : \Term \to \Term
\]
defined by replicating, in a single let-chain, the construction
\begin{align*}
  \mathit{cw} &:= \closeW\,w \\
  \kmax       &:= \mathsf{gRec}\,O\,(s\,\mathit{cw}) \\
  x'          &:= \outK_{\Lstar}\,\kmax \\
  n_T         &:= \text{(sigma-pile fuel in }\kmax,\,
                    \mathsf{fuelMu\_fun}\,\kmax\,\kmax,\,
                    \mathsf{fuelG}\,\kmax\,O\text{)} \\
  \cgFun\,w   &:= \mathsf{cmp}\,
                  (\mathsf{cmp}\,(\mathsf{exfProof}\,D\,\codeFalse)\,\mathsf{cPos})\,
                  (\mathsf{cmp}\,\mathsf{outerWrap}\,\mathsf{cSizeProofDef}),
\end{align*}
where $D$, $\mathsf{cPos}$, $\mathsf{outerWrap}$ are the let-bound
syntactic functions of $\kmax$, $n_T$, $x'$ documented inside the
file. The companion theorem is
\begin{align*}
  \cgFalse :{}
    & (w : \Term) \to {} \\
    & \Deriv\,(\mathsf{eqF}\,(\thmT\,w)\,(\Kcode_{\Lstar}\,(\outK_{\Lstar}\,w))) \to {} \\
    & \Deriv\,(\mathsf{eqF}\,(\thmT\,(\cgFun\,w))\,\codeFalse),
\end{align*}
whose proof is the second $\Sigma$-projection
$\cgFalse\,w\,d := \mathsf{snd}\,(\CGIcoreself\,w\,d)$, with the bridge
$\cgFun\,w \equiv \mathsf{fst}\,(\CGIcoreself\,w\,d)$ holding by
$\mathsf{refl}$ thanks to Agda's let-substitution and record-projection
unfolding.

\section{The subtlety: internal evaluation semantics and the correctness proof}

\subsection{$\cgFun$ is a meta function, not a $\Fun{1}$}

The reader might hope for $\cgFun : \Fun{1}$, i.e.\ an element of the
BRA syntactic algebra $\{s, o, u, C\}$ over $\Fun{2} = \{v, R\}$. This
is provably \emph{not possible} for the present construction. The Term
$\cgFun\,w$ contains $\closeW\,w := \mathsf{substT}\,0\,O\,
(\mathsf{substT}\,1\,O\,w)$, a meta-Agda function that pattern-matches
on the constructor of $w$:
\[
  \closeW\,(\mathsf{var}\,0) = O,
  \quad \closeW\,(\mathsf{var}\,5) = \mathsf{var}\,5.
\]
A BRA $\Fun{1}\,f$ over $\{s, o, u, C\}$ applied to $w$ yields a Term
$\mathsf{ap}_1\,f\,w$ whose structure is fixed by $f$ and which can
only \emph{embed} $w$ at fixed leaf positions; it cannot dispatch on
$w$'s head constructor. The two witnesses
$w \in \{\mathsf{var}\,0,\,\mathsf{var}\,5\}$ already force any
candidate to produce different shapes ($O$ vs.\ $\mathsf{var}\,5$),
which no closed combinator tree achieves. The closeW step is therefore
intrinsically meta-level.

The object-level substitution $\mathsf{sbt} : \Fun{2}$ shipped in
\texttt{BRA4.SbT} does not help: $\mathsf{sbt}$ acts on the Cantor
\emph{encoding} of Terms (firing on $\mathsf{ap}_2\,\mathsf{Pair}\,
(\mathsf{natCode}\,\mathsf{tag\_var})\,(\mathsf{natCode}\,k)$, not on
raw $\mathsf{var}\,k$).

\subsection{Why $\closeW$ is needed in the first place}

The module \texttt{BRA4.StepU2MuCorrect.Construct} that internally
proves the mu-loop run
\[
  \mathsf{runs\_mu} :
    \Deriv\bigl(\mathsf{iter}\,\mathsf{step}\,
        (\mathsf{cfgEV}\,(\mathsf{mcodeMu}\,(\mathsf{mcode1}\,g_F))\,
                          x_{\mathrm{outer}}\,K_0)\,
        (\mathsf{fuel})
      = \mathsf{cfgRT}\,\kmax\,K_0\bigr)
\]
requires three meta-level closedness conditions on its $\kmax$ parameter:
\begin{align*}
  k_{\max\_\mathsf{sub}_0} &: (a : \Term) \to \kmax[0 := a] = \kmax \\
  k_{\max\_\mathsf{sub}_1} &: (a : \Term) \to \kmax[1 := a] = \kmax \\
  k_{\max\_\mathsf{sim}}   &: (a\,b : \Term) \to \kmax[\![0 := a,\,1 := b]\!] = \kmax
\end{align*}
(all are \emph{meta} $\mathsf{Eq}$, not $\Deriv$). These are consumed
by $\mathsf{eqSubst}$-transports inside the
$\mathsf{ruleIndNat}$-with-substituted-motive proof of $\mathsf{runs\_mu}$.

Since $\kmax = \mathsf{gRec}\,O\,(s\,(\closeW\,w))$, structural
induction on $w$ (in \texttt{BRA4.CloseW}) discharges all three
conditions for $\closeW\,w$. Without $\closeW$, the conditions fail at,
e.g., $w := \mathsf{var}\,0$.

\subsection{The internal evaluation correctness chain}

The $\Deriv$-half $\cgFalse$ relies on closing the chain
\[
  \evalU\,(\parse\,\mathsf{gLnameDef})\,n_T
    \stackrel{?}{=} s\,(\outK_{\Lstar}\,\kmax)
\]
through seven internal segments (\texttt{ChainKdef}):
\texttt{stepU\_at\_evC}, \texttt{runs\_mu}, \texttt{stepU\_at\_rtC1},
\texttt{stepU\_at\_evU}, \texttt{stepU\_at\_rtApp2},
\texttt{runs2}\,$(\mathsf{Lift1}\,\outK_{\Lstar})$,
\texttt{stepU\_at\_rtEmpty}, glued by \texttt{iter\_add\_T}'s
fuel-composition lemma. This run is then internalised as the
\textsf{thmT}-fact
\[
  \thmT\,\mathsf{cPos} \stackrel{\Deriv}{=}
    \codeF{\mathsf{runProg}\,g_L\,n_T = s\,x'}
\]
via \texttt{Thm 13} (singulary). The negative leg
$\thmT\,(\mathsf{cmp}\,\mathsf{outerWrap}\,\mathsf{cSize}) =
\mathsf{cNeg}\,\codeF{\mathsf{runProg}\,g_L\,n_T = s\,x'}$ is obtained
from the hypothesis on $w$ by two $\mathsf{thmT\_at\_sb}$
substitutions (\texttt{passK}) followed by $\mathsf{encoded\_mp}$
against \texttt{dLenStarDef}'s size atom. The clash is then a single
$\mathsf{encoded\_mp}\,/\,\mathsf{encoded\_exfalso}$ pair.

\subsection{No consistency hypothesis}

Chaitin's original argument concludes ``$T$ is inconsistent'' (in the
metatheory). The BRA4 formalisation strengthens this: the conclusion
is an \emph{internal} $\Deriv\,(\thmT\,(\cgFun\,w) = \codeFalse)$,
i.e.\ a BRA derivation that the constructed code $\cgFun\,w$ is a
$T$-proof of $0=1$. No $\mathrm{Con}(T)$ premise is consumed, no
\textsf{isNat} integrality predicate is consumed, and the whole
self-referential firing condition is stated in $\num$-raw form (the
subject slot is $\num\,x'$ for an arbitrary Term $x'$, no integrality
required to align with the read-back).

\section{Files}

\begin{itemize}
\item \texttt{BRA4/ChaitinG1CoreNumRaw.agda} --- $\CGIcore$
  ($x$-parametric form, the workhorse).
\item \texttt{BRA4/ChaitinG1CoreNumRawSelf.agda} --- $\CGIcoreself$
  (the $x$-free, self-referential form, one-line specialisation of
  the workhorse).
\item \texttt{BRA4/CgFun.agda} --- $\cgFun : \Term \to \Term$ and
  $\cgFalse$, the $\Sigma$-decomposition.
\item \texttt{BRA4/CloseW.agda} --- the three structural-induction
  lemmas $\mathsf{cl\_w\_sub}_0, \mathsf{cl\_w\_sub}_1,
  \mathsf{cl\_w\_sim}$ enabling the $\Construct$ closure conditions
  at $\kmax(w)$.
\end{itemize}

All files typecheck under \texttt{--safe --without-K --exact-split}
with no holes and no postulates.

\begin{thebibliography}{9}

\bibitem{Chaitin1971}
G.~J. Chaitin, ``Computational complexity and G\"odel's incompleteness theorem,''
\emph{ACM SIGACT News}, no.~9, pp.~11--12, 1971.

\bibitem{Chaitin1974}
G.~J. Chaitin, ``Information-theoretic limitations of formal systems,''
\emph{Journal of the ACM}, vol.~21, no.~3, pp.~403--424, 1974.

\end{thebibliography}

\end{document}
